% engines: lualatex, pdflatex
% expect-warnings: 0
% expect-errors: 0
\documentclass{article}
\usepackage[margin=1cm]{geometry}
\usepackage{../../latex/pgf-go}
\newcommand\case[1]{\par\medskip\noindent\textbf{#1}\par\nopagebreak}
\pagestyle{empty}
\begin{document}

\case{A named goban (name with a colon)}
\begin{goban}[9, name = gob:a]
	\stones[b]{C3, D4, E5}\stones[w]{D3, E4}
	\moves*{F6, G7}
	\marks[b, triangle]{C3}
	\marks[circle]{A1}
	\arrow{B2}{C3}
	\territory[black]{A9, B9}
\end{goban}

\case{One point in the text: D4 \pgfgoview{gob:a}{D4}, E4 \pgfgoview{gob:a}{E4}, A1 (corner, empty) \pgfgoview{gob:a}{A1}, B2 \pgfgoview{gob:a}{B2}}
Text before \pgfgoview{gob:a}{C3}\pgfgoview{gob:a}{D3} and after; a long line to see that the
views do not change the space between lines: x\pgfgoview{gob:a}{E5}x \pgfgoview{gob:a}{F6} \pgfgoview{gob:a}{G7}.

\case{A rectangle: everything by default, then without marks / numbers / territory}
\pgfgoview{gob:a}{A1:G7} \pgfgoview[marks = false]{gob:a}{A1:D4} \pgfgoview[numbers = false]{gob:a}{E5:G7}
\pgfgoview[territory = false, hoshi = false]{gob:a}{A8 to B9}

\case{goban options in a view: no background, no labels, no grid, scale}
\pgfgoview[background = false, label = false]{gob:a}{C3:E5} \pgfgoview[grid = false]{gob:a}{C3:E5}
\pgfgoview[scale = 1.5, raise = 0pt]{gob:a}{D4}

\case{Defaults for all views with \textbackslash pgfgoset\{view\}}
{\pgfgoset{view = {marks = false, background = false}}\pgfgoview{gob:a}{B2:D4}}
\pgfgoview{gob:a}{B2:D4}

\case{resume brings back everything that is asked for (and name keeps it)}
\begin{goban}[9, resume = {gob:a, numbers, marks, territory}, name = b, partial = A1:E5]\stone[w]{B4}\end{goban}
\pgfgoview{b}{C3:D4} \pgfgoview{b}{B4}

\case{The options of the named goban (size, labels, players) are used again; label = false removes the labels}
\begin{goban}[5, name = small, label, players = {white = {stone = {fill = {color = yellow}}}}]\stones[w]{C3}\stones[b]{E5}\end{goban}
\pgfgoview{small}{C3} \pgfgoview{small}{D4:E5} \pgfgoview[label = false]{small}{D4:E5}

\case{A view does not change the memory (resume: C3, D4, E5, D3, E4, F6, G7 and the marks)}
\begin{goban}[9, resume = {gob:a, marks}]\end{goban}

\case{A view keeps the style of the named goban (profile, \textbackslash pgfgoset): wain and roman numbers; a default goban inside wain; view options and \textbackslash pgfgoset\{view\} go on top}
{\pgfgoset{profile = wain}\pgfgoset{move = {label = {format = roman}}}%
\begin{goban}[5, name = styled]\stones[b]{C3}\moves*[w]{D3}\marks[b, triangle]{C3}\end{goban}}
\pgfgoview{styled}{B2:D4} {\pgfgoset{profile = wain}\pgfgoview{gob:a}{C3:E5}}
\pgfgoview[players = {white = {stone = {fill = {color = yellow}}}}]{styled}{B2:D4}
{\pgfgoset{view = {players = {black = {stone = {fill = {color = green}}}}}}\pgfgoview{styled}{B2:D4}}

\case{view/profile: false (initially: the style of the goban), simple, wain, and pgfgo for all views; options of the view go after the profile}
\pgfgoview{styled}{B2:D4} \pgfgoview[profile = simple]{styled}{B2:D4} \pgfgoview[profile = wain]{gob:a}{C3:E5}
{\pgfgoset{view = {profile = pgfgo}}\pgfgoview{styled}{B2:D4} \pgfgoview[stone = {line = {color = red}}]{styled}{B2:D4}}

\case{A sequence keeps the letters of the whole board in a view and in a partial goban (K, O, E-F, J-O; 5, 6, 7)}
\begin{goban}[13, name = seq]
	\stones[b]{B9, C9, D9}\stones[w]{B8, C8, D8}
	\marks[sequence = {format = Alph}]{N10, N9, N8, N7, C2, C1, D1, E1, F1, B10, A10, A9, A8, A7, B7}
\end{goban}
\pgfgoview{seq}{A10} \pgfgoview{seq}{B7} \pgfgoview{seq}{C1:C2} \pgfgoview{seq}{A7:B10}
\begin{goban}[13, partial = A1:F10, scale = 0.5]\marks[sequence]{N10, N9, N8, N7, C2, C1, D1}\end{goban}

\end{document}
